% ============================================================
%  Altar Valley Berms — Figure Registry Report  (Paper 2)
%  Paper 2: Distinguishing breaches from flanks
%  Auto-generated from figures/paper2/figure_registry_concise.txt
%  Date: TBD
% ============================================================
\documentclass[11pt, a4paper]{article}

\usepackage[utf8]{inputenc}
\usepackage[T1]{fontenc}
\usepackage{lmodern}
\usepackage[left=2.5cm, right=2.5cm, top=2.5cm, bottom=2.5cm]{geometry}
\usepackage{graphicx}
\usepackage{tabularx}
\usepackage{booktabs}
\usepackage{caption}
\usepackage{float}
\usepackage{xcolor}
\usepackage{parskip}
\usepackage{microtype}
\usepackage{hyperref}
\hypersetup{colorlinks=true, linkcolor=black, urlcolor=blue}

% ── column types ─────────────────────────────────────────────
\newcolumntype{L}{>{\raggedright\arraybackslash}X}   % wrapping left col

% ── helper: shaded row label ─────────────────────────────────
\newcommand{\rowlbl}[1]{\textbf{\small #1}}

% ── figure + registry table macro ────────────────────────────
%   \figentry{label}{file}{updated}{stats text}{interp text}
\newcommand{\figentry}[5]{%
  \subsection*{#1}
  \begin{figure}[H]
    \centering
    \includegraphics[width=\linewidth]{../figures/paper2/#2}
    \caption*{\textit{File:} \texttt{#2} \quad \textit{Updated:} #3}
  \end{figure}
  \vspace{0.5em}
  \begin{tabularx}{\linewidth}{@{} p{3.2cm} L @{}}
    \toprule
    \rowlbl{Field} & \rowlbl{Content} \\
    \midrule
    \rowlbl{Figure} & #1 \\[2pt]
    \rowlbl{File} & \texttt{\small #2} \\[2pt]
    \rowlbl{Updated} & #3 \\[6pt]
    \rowlbl{Statistical test} & #4 \\[6pt]
    \rowlbl{Interpretation} & #5 \\
    \bottomrule
  \end{tabularx}
  \vspace{3em}
}

% ============================================================
\begin{document}

\begin{center}
  {\LARGE \textbf{Altar Valley Berms — Paper 2}} \\[0.4em]
  {\large Distinguishing Breaches from Flanks} \\[0.3em]
  {\large Figure Registry Report} \\[0.3em]
  {\small Generated TBD}
\end{center}

\vspace{1em}
\hrule
\vspace{1.5em}

This document summarises each registered figure for Paper 2 (breach vs.\ flank
analysis). Add figures below using \verb|\figentry| as they are created and
registered in \texttt{figures/paper2/figure\_registry.txt}.

\vspace{1em}

%% FIG_2_START
\figentry%
  {Figure 2}%
  {fig2.png}%
  {2026-03-10 09:21}%
  {Grouped comparison of berm outcome proportions between Long and Short berm classes, with per-metric 2×2 contingency-table significance tests shown as bracket annotations.}%
  {Outcome proportions differ by length class, and this figure serves as Paper 2 Figure 2 for the breach/flank-focused comparison.}
%% FIG_2_END

%% SI_TABLE_2_START
\begin{table}[H]
  \centering
  \caption{SI Table 2: Berm outcome proportions by length class (threshold = 60\,m). $\Delta$ = P(Long) $-$ P(Short) as a decimal proportion. $p$-values from one-sided Fisher exact tests in direction of the observed difference. Updated 2026-03-02 07:08.}
  \label{tab:si_table2_length_outcomes}
  \small
  \begin{tabular}{lrrrlrl}
    \toprule
    Outcome & P(Long $>$ 60\,m) & P(Short $\leq$ 60\,m) & $\Delta$ & Direction & $p$ (1-sided) & Sig. \\
    \midrule
    Intact & 0.471 & 0.704 & -0.234 & Long < Short & 0.000 & *** \\
    Any failure & 0.529 & 0.296 & 0.234 & Long > Short & 0.000 & *** \\
    Breach & 0.218 & 0.080 & 0.138 & Long > Short & 0.000 & *** \\
    Flank & 0.311 & 0.216 & 0.096 & Long > Short & 0.002 & ** \\
    Vegetation response & 0.487 & 0.464 & 0.023 & Long > Short & 0.284 & ns \\
    Breach | failed & 0.412 & 0.271 & 0.141 & Long > Short & 0.008 & ** \\
    Flank | failed & 0.588 & 0.729 & -0.141 & Long < Short & 0.008 & ** \\
    \bottomrule
  \end{tabular}
\end{table}
%% SI_TABLE_2_END


%% FIG_3_START
\figentry%
  {Figure 3}%
  {fig_texture_fail_type.png}%
  {2026-03-16 05:26}%
  {Top-3 texture × failure type (Intact / Breach / Flank). χ² = 30.82, p = 3.34e-06, Cramér's V = 0.156. 2 significant pairwise contrasts.}%
  {Three-category failure breakdown across top textures. Sandy loam shows highest flanking; clay loam has highest breach share.}
%% FIG_3_END


%% FIG_4_START
\figentry%
  {Figure 4}%
  {fig_landform_fail_type.png}%
  {2026-03-16 05:26}%
  {Landform × failure type (Intact / Breach / Flank). χ² = 14.41, p = 0.00609, Cramér's V = 0.096. 1 significant pairwise contrasts.}%
  {Failure mode breakdown across landforms. Fan terraces show highest flanking rate; flood plains highest breach rate.}
%% FIG_4_END


%% FIG_5_START
\figentry%
  {Figure 5}%
  {fig_length_fail_type.png}%
  {2026-03-16 05:26}%
  {Length × failure type (threshold 50 m). χ² = 46.82, p = 6.8e-11, Cramér's V = 0.246. Fisher exact (breach vs flank × long vs short): OR = 1.95, p = 0.0192.}%
  {Two-panel view of how berm length relates to failure mode. Panel A shows unconditional outcome probabilities; Panel B shows the breach-vs-flank split among failed berms only.}
%% FIG_5_END


%% FIG_6_START
\figentry%
  {Figure 6}%
  {fig_slope_fail_type.png}%
  {2026-03-16 05:26}%
  {Slope × failure type (Intact / Breach / Flank). χ² = 13.79, p = 0.00101, Cramér's V = 0.133. 1 significant pairwise contrasts.}%
  {Failure mode breakdown by slope class. Steeper slopes show higher flanking rate.}
%% FIG_6_END


%% FIG_7_START
\figentry%
  {Figure 7}%
  {fig_soildev_fail_type.png}%
  {2026-03-16 05:26}%
  {Soil development × failure type (Intact / Breach / Flank). χ² = 13.80, p = 0.00101, Cramér's V = 0.133. 1 significant pairwise contrasts.}%
  {Failure mode breakdown by soil development. B-horizon soils show higher flanking rate.}
%% FIG_7_END


%% FIG_8_START
\figentry%
  {Figure 8}%
  {fig_stacked_fail_by_landform.png}%
  {2026-03-16 05:26}%
  {100 % stacked bar chart showing relative Intact / Breach / Flank composition within each landform.}%
  {Proportional failure-type composition by landform. Flood plains have the highest breach share; fan terraces the highest flank share.}
%% FIG_8_END


%% FIG_9_START
\figentry%
  {Figure 9}%
  {fig_angle_fail_type.png}%
  {2026-03-16 05:26}%
  {Absolute berm angle by failure type. Kruskal-Wallis p = 5.14e-07.}%
  {Box + strip plot of berm directionality by Fail_Type. Tests whether berms farther from perpendicular are more prone to a specific failure mode.}
%% FIG_9_END


%% FIG_10_START
\figentry%
  {Figure 10}%
  {fig_flowaccum_fail_type.png}%
  {2026-03-16 05:26}%
  {Log flow accumulation by failure type. Kruskal-Wallis p = 4.86e-10.}%
  {Box + strip plot of log-transformed upstream flow accumulation by Fail_Type. Tests whether berms with larger contributing catchments are more prone to breach or flank.}
%% FIG_10_END


%% FIG_11_START
\figentry%
  {Figure 11}%
  {fig_savi_fail_type.png}%
  {2026-03-16 05:26}%
  {SAVI vegetation indices by failure type. KW p: downstream = 0.377, upstream = 0.876, veg effect = 0.911.}%
  {Three-panel view of downstream SAVI, upstream SAVI, and normalised vegetation effect by Fail_Type.}
%% FIG_11_END


%% FIG_12_START
\figentry%
  {Figure 12}%
  {fig_roaddist_fail_type.png}%
  {2026-03-16 05:26}%
  {Road distance by failure type. Kruskal-Wallis p = 0.00901.}%
  {Box + strip plot of distance to nearest road by Fail_Type. Tests whether road proximity influences failure mode.}
%% FIG_12_END


%% FIG_13_START
\figentry%
  {Figure 13}%
  {fig_highclay_fail_type.png}%
  {2026-03-16 05:26}%
  {High vs Low clay × failure type. χ² = 18.49, p = 9.66e-05, Cramér's V = 0.154.}%
  {Failure mode breakdown by clay content. Higher clay soils show relatively more breaching and less flanking.}
%% FIG_13_END

\end{document}
